\chapter{2026年北京卷}

\section{单选题}

\begin{problem}
已知集合 \( M=\left\{ x\mid -1<x<3 \right\} \)，\( N=\left\{ x\mid x\ge 2 \right\} \)，则 \( M\cup N=\)（\quad）

\choices
    {\( \left( 2,3 \right) \)}
    {\( \left( -1,+\infty \right) \)}
    {\( \left[ 2,3 \right) \)}
    {\( \left( -1,2 \right] \)}

\end{problem}

\begin{answer}
B
\end{answer}

\begin{solution}
因为 \( M=\left( -1,3 \right) \)，\( N=\left[ 2,+\infty \right) \)，所以 \( M\cup N=\left( -1,+\infty \right) \). 故选 B.
\end{solution}

\begin{problem}
已知 \( z_1=3-2\i \)，\( z_2=-1+4\i \)，则 \( \left| z_1+z_2 \right|=\)（\quad）

\choices
    {\( 2\sqrt{2} \)}
    {\( \i \)}
    {2}
    {8}

\end{problem}

\begin{answer}
A
\end{answer}

\begin{solution}
由 \( z_1+z_2=\left( 3-2\i \right)+\left( -1+4\i \right)=2+2\i \)，得 \( \left| z_1+z_2 \right|=\left| 2+2\i \right|=\sqrt{2^2+2^2}=2\sqrt{2} \). 故选 A.
\end{solution}

\begin{problem}
已知双曲线 \( C:\frac{x^2}{a^2}-\frac{y^2}{4}=1 \)（\( a>0 \)）的渐近线方程为 \( y=\pm \frac{2}{3}x \)，则 \( a \) 的值为（\quad）

\choices
    {2}
    {3}
    {4}
    {9}

\end{problem}

\begin{answer}
B
\end{answer}

\begin{solution}
双曲线 \( \frac{x^2}{a^2}-\frac{y^2}{4}=1 \) 的渐近线方程为 \( y=\pm \frac{2}{a}x \). 由 \( \frac{2}{a}=\frac{2}{3} \)，得 \( a=3 \). 故选 B.
\end{solution}

\begin{problem}
已知 \( \left( a-\sqrt{x} \right)^7 \) 的展开式中的 \( x^2 \) 的系数是 280，则 \( a=\)（\quad）

\choices
    {2}
    {\(-2\)}
    {1}
    {\(-1\)}

\end{problem}

\begin{answer}
A
\end{answer}

\begin{solution}
在 \( \left( a-\sqrt{x} \right)^7 \) 的展开式中，\( x^2 \) 项来自 \( \mathrm C_7^4 a^3\left( -\sqrt{x} \right)^4=35a^3x^2 \). 所以 \( 35a^3=280 \)，即 \( a^3=8 \). 因为 \( a=2 \)，故选 A.
\end{solution}

\begin{problem}
下列函数 \( f(x) \) 是奇函数且在定义域上单调递增的是（\quad）

\choices
    {\( f(x)=\frac{1}{x^2+1} \)}
    {\( f(x)=\sin x \)}
    {\( f(x)=2^{-x}-2^x \)}
    {\( f(x)=\ln \frac{5+x}{5-x} \)}

\end{problem}

\begin{answer}
D
\end{answer}

\begin{solution}
选项 A 不是奇函数. 选项 B 虽是奇函数，但在定义域上不单调递增. 选项 C 满足 \( f\left( -x \right)=-f\left( x \right) \)，但 \( f'\left( x \right)=-\ln 2\left( 2^{-x}+2^x \right)<0 \)，它单调递减.
对选项 D，定义域为 \( \left( -5,5 \right) \). 因为 \( f\left( -x \right)=\ln \frac{5-x}{5+x}=-\ln \frac{5+x}{5-x}=-f\left( x \right) \)，所以它是奇函数. 又 \( f'\left( x \right)=\frac{1}{5+x}+\frac{1}{5-x}=\frac{10}{25-x^2}>0 \)，故它在定义域上单调递增. 故选 D.
\end{solution}

\begin{problem}
已知向量 \( \bs{a} \)，\( \bs{b} \) 满足 \( \left| \bs{a}-\bs{b} \right|=2 \)，\( \bs{a}=\left( 2,0 \right) \)，则 \( \left| \bs{b} \right| \) 的最大值为（\quad）

\choices
    {1}
    {2}
    {3}
    {4}

\end{problem}

\begin{answer}
D
\end{answer}

\begin{solution}
由 \( \left| \bs{a}-\bs{b} \right|=2 \) 知，向量 \( \bs{b} \) 的终点在以 \( \left( 2,0 \right) \) 为圆心，半径为 2 的圆上. 设 \(O\) 为原点，则 \( \left| \bs{b} \right|\le \left| \bs{a} \right|+\left| \bs{a}-\bs{b} \right|=2+2=4 \). 当 \( \bs{b} \) 与 \( \bs{a} \) 同向，且 \( \left| \bs{b} \right|=4 \) 时取等号. 故选 D.
\end{solution}

\begin{problem}
\( \left\{ a_n \right\} \)，\( \left\{ b_n \right\} \) 是无穷数列，则“存在常数 \(M\)，使 \( a_n\le M\le b_n \)（\(n=1\)，\(2\)，\(3\)，\(\cdots\)）”是“\( a_n\le b_n \)（\(n=1\)，\(2\)，\(3\)，\(\cdots\)）”的（\quad）

\choices
    {充分不必要条件}
    {必要不充分条件}
    {充要条件}
    {既不充分也不必要条件}

\end{problem}

\begin{answer}
A
\end{answer}

\begin{solution}
若存在常数 \(M\)，使 \( a_n\le M\le b_n \) 对任意 \( n \) 都成立，则必有 \( a_n\le b_n \). 所以前者是后者的充分条件.
它不是必要条件. 例如取 \( a_n=n \)，\( b_n=n+1 \)，则 \( a_n\le b_n \) 对任意 \( n \) 都成立，但不存在常数 \(M\) 使 \( n\le M\le n+1 \) 对任意 \( n \) 都成立.
所以前者是后者的充分不必要条件. 故选 A.
\end{solution}

\begin{problem}
\( f(x)=\sin\left( x+\varphi \right) \)（\( 0<\varphi<2\pi \)），将 \( f(x) \) 向 \( x \) 轴正方向平移 \( 3\varphi \) 个单位，得到的函数图像与 \( f(x) \) 图像关于 \( x \) 轴对称，则 \( \varphi \) 的取值个数为（\quad）

\choices
    {1}
    {2}
    {3}
    {4}

\end{problem}

\begin{answer}
C
\end{answer}

\begin{solution}
将 \( f(x) \) 向右平移 \( 3\varphi \) 个单位后，得到 \( g(x)=f\left( x-3\varphi \right)=\sin\left( x-2\varphi \right) \).
由题意，\( g(x)=-f(x) \) 对任意 \( x \) 都成立，所以 \( \sin\left( x-2\varphi \right)+\sin\left( x+\varphi \right)=0 \)，即 \( 2\sin\left( x-\frac{\varphi}{2} \right)\cos\frac{3\varphi}{2}=0 \).
对任意 \( x \) 恒成立，只能有 \( \cos\frac{3\varphi}{2}=0 \). 所以 \( \frac{3\varphi}{2}=\frac{\pi}{2}+k\pi \)，从而 \( \varphi=\frac{\pi}{3}+\frac{2k\pi}{3} \).
结合 \( 0<\varphi<2\pi \)，得 \(\varphi=\frac{\pi}{3}\)，\(\pi\)，\(\frac{5\pi}{3}\)，共有 3 个. 故选 C.
\end{solution}

\begin{problem}
学校组织高一，高二学生参观甲，乙两地博物馆，每位学生可自主选择一处前往. 已知高一学生总人数多于高二学生总人数，前往甲地的全体学生总数多于前往乙地的全体学生总数，则（\quad）

\choices
    {去甲地的高一学生人数多于去乙地的高一学生人数}
    {去甲地的高一学生人数多于去乙地的高二学生人数}
    {去甲地的高一学生人数不多于去乙地的高二学生人数}
    {去乙地的高二学生人数不少于去甲地的高二学生人数}

\end{problem}

\begin{answer}
B
\end{answer}

\begin{solution}
设去甲地的高一，高二学生人数分别为 \(x\)，\(u\)，去乙地的高一，高二学生人数分别为 \(y\)，\(v\). 由题意得 \( x+y>u+v \) 且 \( x+u>y+v \).
两式相加，得 \( 2x>2v \)，所以 \( x>v \). 这正是“去甲地的高一学生人数多于去乙地的高二学生人数”.
因此选项 B 一定成立. 故选 B.
\end{solution}

\begin{problem}
摇杆机械装置，如图，\(A\)，\(B\) 为定点，\(C\)，\(D\) 是动点，\( AD=1 \)，\( CD=3 \)，\( BC=\frac{5}{2} \)，\( AB=4 \)，则 \( \cos \angle ABC \) 的取值范围（\quad）

\bitmapfigure[max width=0.42\linewidth,max totalheight=3.9cm]{img/2026/beijing/q10_fig1.png}

\choices
    {\( \left[ \frac{5}{16},\frac{73}{80} \right] \)}
    {\( \left[ -\frac{5}{16},\frac{73}{80} \right] \)}
    {\( \left[ \frac{5}{16},1 \right] \)}
    {\( \left[ \frac{5}{16},\frac{73}{80} \right) \)}

\end{problem}

\begin{answer}
A
\end{answer}

\begin{solution}
在三角形 \(ADC\) 中，已知 \( AD=1 \)，\( CD=3 \)，所以由三角形两边之和大于第三边，两边之差小于第三边，得 \( 2\le AC\le 4 \). 两个端点都可取到，此时 \(A\)，\(D\)，\(C\) 三点共线.
在三角形 \(ABC\) 中，由余弦定理得 \( \cos \angle ABC=\frac{AB^2+BC^2-AC^2}{2\cdot AB\cdot BC}=\frac{16+\frac{25}{4}-AC^2}{20}=\frac{89}{80}-\frac{AC^2}{20} \).
它随 \(AC\) 的增大而减小. 当 \( AC=2 \) 时，\( \cos \angle ABC=\frac{89}{80}-\frac{4}{20}=\frac{73}{80} \). 当 \( AC=4 \) 时，\( \cos \angle ABC=\frac{89}{80}-\frac{16}{20}=\frac{5}{16} \).
所以 \( \cos \angle ABC \) 的取值范围为 \( \left[ \frac{5}{16},\frac{73}{80} \right] \). 故选 A.
\end{solution}
\section{填空题}

\begin{problem}
已知直线 \( ax+y=0 \) 与圆 \( \left( x-2 \right)^2+\left( y-2 \right)^2=4 \) 相切，则 \( a=\fillinblank{} \).

\end{problem}

\begin{answer}
0
\end{answer}

\begin{solution}
圆心为 \( \left( 2,2 \right) \)，半径为 2. 直线 \( ax+y=0 \) 到圆心的距离为 \( \frac{\left| 2a+2 \right|}{\sqrt{a^2+1}} \). 由相切得 \( \frac{\left| 2a+2 \right|}{\sqrt{a^2+1}}=2 \). 两边平方，得 \( \left( a+1 \right)^2=a^2+1 \).
所以 \( 2a=0 \)，即 \( a=0 \).
\end{solution}

\begin{problem}
已知等差数列 \( \left\{ a_n \right\} \) 的前 \( n \) 项和为 \( S_n \)，且 \( S_6=6a_6+30 \)，则公差 \( d=\fillinblank{} \)，若 \( S_n\le S_5 \) 恒成立，则符合条件的 \( a_1 \) 的一个取值为 \(\fillinblank{} \).

\end{problem}

\begin{answer}
\(-2\)，\(8\)（第二空答案不唯一）
\end{answer}

\begin{solution}
由 \( S_6=\frac{6}{2}\left[ 2a_1+5d \right]=6a_6+30=6\left( a_1+5d \right)+30 \) 得 \( 6a_1+15d=6a_1+30d+30 \)，所以 \( d=-2 \).
于是 \( a_n=a_1+\left( n-1 \right)d=a_1-2n+2 \)，且 \( S_n=\frac{n}{2}\left[ 2a_1+\left( n-1 \right)\left( -2 \right) \right]=n\left( a_1-n+1 \right) \).
由 \( S_n\le S_5 \) 恒成立，知 \(S_5\) 是该数列前 \( n \) 项和的最大值. 对等差数列的前 \( n \) 项和，等价于 \( a_5\ge 0 \) 且 \( a_6\le 0 \). 即 \( a_1-8\ge 0 \) 且 \( a_1-10\le 0 \)，所以 \( 8\le a_1\le 10 \). 例如可取 \( a_1=8 \).
\end{solution}

\begin{problem}
音高 \( y \)（单位：mel）与频率 \( f \)（单位：Hz）满足 \( y=\lg \left( 1+\frac{f}{700} \right) \)，若 \( \lg 2\le y<3\lg 2 \)，则 \( f \) 的取值范围为 \(\fillinblank{} \).

\end{problem}

\begin{answer}
\( \left[ 700,4900 \right) \)
\end{answer}

\begin{solution}
由 \( \lg 2\le \lg \left( 1+\frac{f}{700} \right)<3\lg 2=\lg 8 \)
得 \( 2\le 1+\frac{f}{700}<8 \)，所以 \( 1\le \frac{f}{700}<7 \)，从而 \( 700\le f<4900 \). 因此 \( f \) 的取值范围为 \( \left[ 700,4900 \right) \).
\end{solution}

\begin{problem}
已知三棱锥 \( A-BCD \)，\( AD=AB=AC=2\sqrt{2} \)，\( BD=BC=2 \)，\( DC=2\sqrt{3} \)，则它的底面 \(BCD\) 的面积为 \(\fillinblank{} \)，体积为 \(\fillinblank{} \).

\bitmapfigure[max width=0.34\linewidth,max totalheight=4.2cm]{img/2026/beijing/q14_fig1.png}

\end{problem}

\begin{answer}
\(\sqrt{3}\)，\(\frac{2\sqrt{3}}{3}\)
\end{answer}

\begin{solution}
在三角形 \(BCD\) 中，\( BC=BD=2 \)，\( CD=2\sqrt{3} \). 设底边 \(CD\) 上的高为 \( h \)，则 \( h=\sqrt{2^2-\left( \sqrt{3} \right)^2}=1 \)，所以底面面积 \( S_{BCD}=\frac{1}{2}\cdot 2\sqrt{3}\cdot 1=\sqrt{3} \).
又因为 \( AB=AC=AD \)，所以点 \(A\) 在底面 \(BCD\) 上的射影是三角形 \(BCD\) 的外心. 设外接圆半径为 \(R\)，则 \( R=\frac{BC\cdot BD\cdot CD}{4S_{BCD}}=\frac{2\cdot 2\cdot 2\sqrt{3}}{4\sqrt{3}}=2 \).
设三棱锥高为 \(H\)，则 \( H=\sqrt{AB^2-R^2}=\sqrt{\left( 2\sqrt{2} \right)^2-2^2}=2 \). 所以体积 \( V=\frac{1}{3}S_{BCD}H=\frac{1}{3}\cdot \sqrt{3}\cdot 2=\frac{2\sqrt{3}}{3} \).
故两空分别为 \( \sqrt{3} \) 和 \( \frac{2\sqrt{3}}{3} \).
\end{solution}

\begin{problem}
已知 \( f(x)=\left| x^2-\cos x-c^2 \right| \)，给出下列四个结论：

\begin{circlelist}
\item \( f(x) \) 在 \( \left( -1,1 \right] \) 上有最小值和最大值
\item \( c=0 \)，\( f(x)=1 \) 有 3 个解
\item \( c=1 \)，\( x\in \left( -1,2 \right] \) 时，\( f(x) \) 有最大值
\item \( c>0 \)，\( f(x) \) 与 \( y=c \) 有 4 个交点
\end{circlelist}

其中正确结论的序号是 \(\fillinblank{} \).

\end{problem}

\begin{answer}
\(\circled{1}\)，\(\circled{2}\)，\(\circled{3}\)，\(\circled{4}\)
\end{answer}

\begin{solution}
设 \( g(x)=x^2-\cos x \)，则 \( f(x)=\left| g(x)-c^2 \right| \). 因为 \( g(x) \) 是偶函数，所以 \( f(x) \) 也是偶函数.

\circled{1} 由于 \( f(x) \) 为偶函数，\( f(x) \) 在 \( \left( -1,1 \right] \) 上的最值与它在 \( \left[ 0,1 \right] \) 上的最值相同. 而 \( f(x) \) 在 \( \left[ 0,1 \right] \) 上连续，所以一定有最小值和最大值. 因此 \circled{1} 正确.

\circled{2} 当 \( c=0 \) 时，方程 \( f(x)=1 \) 化为 \( \left| x^2-\cos x \right|=1 \)，即 \( x^2-\cos x=1 \) 或 \( x^2-\cos x=-1 \).
第二个方程等价于 \( x^2=\cos x-1 \). 因为左边非负，右边非正，所以只能有 \( x=0 \).
对第一个方程，设 \( h(x)=x^2-\cos x-1 \). 它是偶函数，且当 \( x>0 \) 时，\( h'\left( x \right)=2x+\sin x>0 \)，
所以 \( h(x) \) 在 \( \left( 0,+\infty \right) \) 上单调递增. 又 \( h(1)=-\cos 1<0 \)，\( h\left( \sqrt{2} \right)=1-\cos \sqrt{2}>0 \)，故有且仅有一个正根，从而也有且仅有一个负根. 所以共有 3 个解，\circled{2} 正确.

\circled{3} 当 \( c=1 \) 时，\( f(x)=\left| x^2-\cos x-1 \right| \). 若 \( x\in \left( -1,1 \right] \)，则 \( f(x)=1+\cos x-x^2\le 2 \)，而 \( f(2)=\left| 4-\cos 2-1 \right|=3-\cos 2>3>2 \). 又当 \( x\in \left[ 1,2 \right] \) 时，\( \left( x^2-\cos x-1 \right)'=2x+\sin x>0 \)，
所以 \( f(x) \) 在 \( \left[ 1,2 \right] \) 上递增，最大值在 \( x=2 \) 处取得. 因此 \circled{3} 正确.

\circled{4} 当 \( c>0 \) 时，方程 \( f(x)=c \) 化为 \( x^2-\cos x=c^2+c \) 或 \( x^2-\cos x=c^2-c \). 注意到 \( c^2-c=\left( c-\frac{1}{2} \right)^2-\frac{1}{4}>-1 \)，
而 \( g(x)=x^2-\cos x \) 是偶函数，且在 \( \left[ 0,+\infty \right) \) 上单调递增，在 \( \left( -\infty,0 \right] \) 上单调递减，最小值为 \( g(0)=-1 \). 所以只要右端大于 \( -1 \)，对应方程就恰有 2 个解. 于是上面两个方程各有 2 个解，共有 4 个交点. 因此 \circled{4} 正确.

综上，正确结论是 \(\circled{1}\)，\(\circled{2}\)，\(\circled{3}\)，\(\circled{4}\).
\end{solution}
\section{解答题}

\begin{problem}
已知函数 \( f(x)=2\sin \omega x\cos \varphi+2\cos \omega x\sin \varphi \)，\( \omega>0 \)，\( \left| \varphi \right|<\frac{\pi}{2} \)，\( f(x) \) 最小正周期为 \( \pi \)，且 \( f(0)>0 \)，\( f\left( \frac{\pi}{4} \right)=1 \).

\begin{enumerate}
\item 求 \( \omega \)，\( \varphi \) 的值；
\item 求 \( f(x) \) 的单调递减区间.
\end{enumerate}

\end{problem}

\begin{solution}
由两角和公式，\( f(x)=2\sin\left( \omega x+\varphi \right) \). 对于 \( y=2\sin\left( \omega x+\varphi \right) \)，最小正周期为 \( \frac{2\pi}{\omega} \). 由题意得 \( \frac{2\pi}{\omega}=\pi \)，所以 \( \omega=2 \). 于是 \( f(x)=2\sin\left( 2x+\varphi \right) \). 由 \( f(0)>0 \) 得 \( 2\sin \varphi>0 \). 又 \( \left| \varphi \right|<\frac{\pi}{2} \)，所以 \( 0<\varphi<\frac{\pi}{2} \). 再由 \( f\left( \frac{\pi}{4} \right)=1 \)，得 \( 2\sin\left( \frac{\pi}{2}+\varphi \right)=1 \)，即 \( 2\cos \varphi=1 \). 所以 \( \cos \varphi=\frac{1}{2} \). 结合 \( 0<\varphi<\frac{\pi}{2} \)，得 \( \varphi=\frac{\pi}{3} \). 因此 \( \omega=2 \)，\( \varphi=\frac{\pi}{3} \).
又 \( f'(x)=4\cos\left( 2x+\frac{\pi}{3} \right) \). 函数单调递减时，\( f'(x)<0 \)，所以 \( \cos\left( 2x+\frac{\pi}{3} \right)<0 \)，即 \( \frac{\pi}{2}+2k\pi<2x+\frac{\pi}{3}<\frac{3\pi}{2}+2k\pi \)，从而 \( \frac{\pi}{12}+k\pi<x<\frac{7\pi}{12}+k\pi \).
其中 \( k\in\Z \). 所以 \( f(x) \) 的单调递减区间为 \( \left( \frac{\pi}{12}+k\pi,\frac{7\pi}{12}+k\pi \right) \)（\( k\in\Z \)）.
\end{solution}

\begin{problem}
现从全校学生中随机抽取 200 人统计数学成绩，成绩分组及对应人数如下：

\begin{center}
\renewcommand{\arraystretch}{1.35}
\begin{tabular}{|c|c|c|c|c|c|}
\hline
成绩分组 & {\([81,94)\)} & {\([94,107)\)} & {\([107,120)\)} & {\([120,135)\)} & {\([135,150]\)} \\
\hline
人数 & 40 & 60 & 60 & 32 & 8 \\
\hline
\end{tabular}
\end{center}

以频率估计概率，完成下列问题：

\begin{enumerate}
\item 求数学成绩低于 120 分的概率；
\item 从学校随机抽取 4 人，求 2 人不低于 120 且 2 人小于 94 的概率；
\item 每组数据取左端，中间，右端，比较 \( s_{\text{左}}^2 \)，\( s_{\text{中}}^2 \)，\( s_{\text{右}}^2 \) 的大小关系.
\end{enumerate}

\end{problem}

\begin{solution}
（1）低于 120 分的人数为 \( 40+60+60=160 \)，所以所求概率为 \( \frac{160}{200}=\frac{4}{5} \).

（2）由频率估计概率，不低于 120 分的概率为 \( \frac{32+8}{200}=\frac{1}{5} \)，小于 94 分的概率也为 \( \frac{40}{200}=\frac{1}{5} \).
从学校随机抽取 4 人时，把每次抽取视为相互独立，且每人落入“不低于 120 分”“小于 94 分”“其余分数段”这三类的概率分别为 \( \frac{1}{5} \)，\( \frac{1}{5} \)，\( \frac{3}{5} \).
恰有 2 人不低于 120 分且 2 人小于 94 分，就是 4 次独立抽取中，出现 2 次第一类，2 次第二类，0 次第三类，所以所求概率为 \( \frac{4!}{2!2!}\left( \frac{1}{5} \right)^2\left( \frac{1}{5} \right)^2=\frac{6}{625} \).

（3）设每个组分别取左端点时，对应的代表值随机变量为 \(X\). 记
\(
\varepsilon=
\begin{cases}
0, & \text{前 3 组}, \\
1, & \text{后 2 组}.
\end{cases}
\)
则每组取中点，右端点时，对应代表值可分别写成 \( X+6.5+\varepsilon \) 和 \( X+13+2\varepsilon \).
因此 \( s_{\text{左}}^2=D\left( X \right) \)，\( s_{\text{中}}^2=D\left( X+\varepsilon \right) \)，\( s_{\text{右}}^2=D\left( X+2\varepsilon \right) \).
又因为 \( D\left( X+t\varepsilon \right)=D\left( X \right)+t^2D\left( \varepsilon \right)+2t\operatorname{Cov}\left( X,\varepsilon \right) \)，而后 2 组对应的分数明显大于前 3 组，所以 \( \operatorname{Cov}\left( X,\varepsilon \right)>0 \)，且 \( D\left( \varepsilon \right)>0 \). 于是 \( s_{\text{中}}^2-s_{\text{左}}^2=D\left( \varepsilon \right)+2\operatorname{Cov}\left( X,\varepsilon \right)>0 \)，且 \( s_{\text{右}}^2-s_{\text{中}}^2=3D\left( \varepsilon \right)+2\operatorname{Cov}\left( X,\varepsilon \right)>0 \)，所以 \( s_{\text{左}}^2<s_{\text{中}}^2<s_{\text{右}}^2 \).
\end{solution}

\begin{problem}
已知直三棱柱 \( ABC-A_1B_1C_1 \)，\( \angle BAC=90^\circ \)，\( AB=AC=2 \)，\( BB_1=\sqrt{2} \)，\(E\)，\(D\) 分别为 \(A_1B_1\)，\(AC\) 的中点.

\begin{enumerate}
\item 证明：\( DE // \text{平面}\, BB_1C_1C \)；
\item 点 \(P\) 在平面 \(A_1B_1C_1\) 内，且 \( EP//B_1C_1 \)，再从条件\(\circled{1}\)，条件\(\circled{2}\)，条件\(\circled{3}\)这三个条件中选择一个作为已知，使得 \(P\) 唯一确定，求平面 \(PAD\) 与平面 \(PDE\) 的夹角的余弦值.

条件\(\circled{1}\)：\( PA=PD \)；

条件\(\circled{2}\)：\( PA\perp BC \)；

条件\(\circled{3}\)：\( BB_1 // \text{平面}\, PDE \).

\end{enumerate}

\bitmapfigure[max width=0.5\linewidth,max totalheight=4.4cm]{img/2026/beijing/q18_fig1.png}

\end{problem}

\begin{solution}
以 \(A\) 为原点，建立空间直角坐标系，取 \( A\left( 0,0,0 \right) \)，\( B\left( 2,0,0 \right) \)，\( C\left( 0,2,0 \right) \). 则 \( A_1\left( 0,0,\sqrt{2} \right) \)，\( B_1\left( 2,0,\sqrt{2} \right) \)，\( C_1\left( 0,2,\sqrt{2} \right) \). 由中点条件，\( D\left( 0,1,0 \right) \)，\( E\left( 1,0,\sqrt{2} \right) \).

（1）有 \( \overrightarrow{DE}=\left( 1,-1,\sqrt{2} \right) \)，而 \( \overrightarrow{BC}=\left( -2,2,0 \right) \)，\( \overrightarrow{BB_1}=\left( 0,0,\sqrt{2} \right) \)，且 \( \overrightarrow{DE}=-\frac{1}{2}\overrightarrow{BC}+\overrightarrow{BB_1} \).
因为 \( \overrightarrow{BC} \) 与 \( \overrightarrow{BB_1} \) 都在平面 \(BB_1C_1C\) 内，故 \( \overrightarrow{DE} \) 平行于该平面. 所以 \( DE // \text{平面}\, BB_1C_1C \).

（2）选条件 \(\circled{2}\)：\( PA\perp BC \).
因为 \(P\) 在平面 \(A_1B_1C_1\) 内，且 \( EP//B_1C_1 \)，所以可设 \( P\left( 1-2t,2t,\sqrt{2} \right) \). 又 \( \overrightarrow{BC}=\left( -2,2,0 \right) \). 由 \( PA\perp BC \)，得 \( \overrightarrow{PA}\cdot \overrightarrow{BC}=0 \)，即 \( \left( -1+2t,-2t,-\sqrt{2} \right)\cdot \left( -2,2,0 \right)=0 \). 化简得 \( t=\frac{1}{4} \)，所以 \( P\left( \frac{1}{2},\frac{1}{2},\sqrt{2} \right) \).
从而 \(P\) 被唯一确定.
取 \( \overrightarrow{AD}=\left( 0,1,0 \right) \) 以及 \( \overrightarrow{AP}=\left( \frac{1}{2},\frac{1}{2},\sqrt{2} \right) \)，则平面 \(PAD\) 的一个法向量可取为 \( \bs{n}_1=\overrightarrow{AD}\times \overrightarrow{AP}=\left( \sqrt{2},0,-\frac{1}{2} \right) \).
又 \( \overrightarrow{DE}=\left( 1,-1,\sqrt{2} \right) \)，\( \overrightarrow{DP}=\left( \frac{1}{2},-\frac{1}{2},\sqrt{2} \right) \)，则平面 \(PDE\) 的一个法向量可取为 \( \bs{n}_2=\overrightarrow{DE}\times \overrightarrow{DP}=\left( -\frac{\sqrt{2}}{2},-\frac{\sqrt{2}}{2},0 \right) \). 所以两平面的夹角余弦为 \( \frac{\left| \bs{n}_1\cdot \bs{n}_2 \right|}{\left| \bs{n}_1 \right|\left| \bs{n}_2 \right|}=\frac{1}{\frac{3}{2}\cdot 1}=\frac{2}{3} \).
故所求夹角的余弦值为 \( \frac{2}{3} \).
\end{solution}

\begin{problem}
已知椭圆 \( E:\frac{x^2}{a^2}+\frac{y^2}{b^2}=1 \)（\( a>b>0 \)）的一个顶点是 \( \left( 2,0 \right) \)，离心率为 \( \frac{1}{2} \).

\begin{enumerate}
\item 求 \(E\) 的方程；
\item 过点 \( A\left( 1,1 \right) \)，斜率为 \( k \)（\( k\neq \pm 1 \)）的直线交椭圆 \(E\) 于 \(B\)，\(C\) 两点，\(B\) 关于 \( y=x \) 的对称点为 \(D\)，\(DC\) 交 \( y=x \) 于 \(Q\)，若 \( \left| S_{\triangle ABQ}-S_{\triangle ACQ} \right|=\frac{5}{8} \)，求 \( k \).
\end{enumerate}

\end{problem}

\begin{solution}
（1）由顶点 \( \left( 2,0 \right) \) 可知 \( a=2 \). 又离心率 \( e=\frac{c}{a}=\frac{1}{2} \)，所以 \( c=1 \). 由 \( c^2=a^2-b^2 \)，得 \( b^2=4-1=3 \).
故椭圆方程为 \( \frac{x^2}{4}+\frac{y^2}{3}=1 \).

（2）过点 \( A\left( 1,1 \right) \) 且斜率为 \( k \) 的直线方程为 \( y=kx-k+1 \).
设 \( B\left( x_1,y_1 \right) \)，\( C\left( x_2,y_2 \right) \). 则 \( y_1=kx_1-k+1 \)，\( y_2=kx_2-k+1 \).
代入椭圆方程，得 \( \frac{x^2}{4}+\frac{\left( kx-k+1 \right)^2}{3}=1 \)，化简为 \( \left( 3+4k^2 \right)x^2+8k\left( 1-k \right)x+4\left( 1-k \right)^2-12=0 \). 于是 \( x_1+x_2=-\frac{8k\left( 1-k \right)}{3+4k^2} \)，\( x_1x_2=\frac{4\left( 1-k \right)^2-12}{3+4k^2} \).
令 \( \alpha=x_1-1 \)，\( \beta=x_2-1 \).
则 \( \alpha\beta=x_1x_2-\left( x_1+x_2 \right)+1=-\frac{5}{3+4k^2} \).
因为 \( \frac{1^2}{4}+\frac{1^2}{3}=\frac{7}{12}<1 \)，
所以点 \( A\left( 1,1 \right) \) 在椭圆内部，故 \(A\) 在线段 \(BC\) 上，从而 \( \alpha\beta<0 \).
由 \(B\) 关于 \( y=x \) 的对称点为 \(D\)，得 \( D\left( y_1,x_1 \right)=\left( 1+k\alpha,1+\alpha \right) \)，而 \( C=\left( 1+\beta,1+k\beta \right) \).
设 \( Q=\left( t,t \right) \). 因为 \( Q\in DC \)，联立直线 \(DC\) 与 \( y=x \) 可得 \( t=1+\frac{\left( 1+k \right)\alpha\beta}{\alpha+\beta} \).
所以 \( AQ=\sqrt{2}\left| \frac{\left( 1+k \right)\alpha\beta}{\alpha+\beta} \right| \).
又因为 \(A\)，\(Q\) 都在直线 \( y=x \) 上，所以 \( S_{\triangle ABQ}=\frac{1}{2}AQ\cdot d\left( B,y=x \right) \)，\( S_{\triangle ACQ}=\frac{1}{2}AQ\cdot d\left( C,y=x \right) \). 其中 \( d\left( B,y=x \right)=\frac{\left| y_1-x_1 \right|}{\sqrt{2}}=\frac{\left| k-1 \right|\left| \alpha \right|}{\sqrt{2}} \)，\( d\left( C,y=x \right)=\frac{\left| y_2-x_2 \right|}{\sqrt{2}}=\frac{\left| k-1 \right|\left| \beta \right|}{\sqrt{2}} \).
因为 \( \alpha\beta<0 \)，所以 \( \left| \left| \alpha \right|-\left| \beta \right| \right|=\left| \alpha+\beta \right| \). 因此 \( \left| S_{\triangle ABQ}-S_{\triangle ACQ} \right|=\frac{1}{2}\cdot \sqrt{2}\left| \frac{\left( 1+k \right)\alpha\beta}{\alpha+\beta} \right|\cdot \frac{\left| k-1 \right|\left| \alpha+\beta \right|}{\sqrt{2}}=\frac{\left| 1-k^2 \right|\left| \alpha\beta \right|}{2} \).
即 \( \left| S_{\triangle ABQ}-S_{\triangle ACQ} \right|=\frac{5\left| 1-k^2 \right|}{2\left( 3+4k^2 \right)} \). 由题意得 \( \frac{5\left| 1-k^2 \right|}{2\left( 3+4k^2 \right)}=\frac{5}{8} \)，化简得 \( 4\left| 1-k^2 \right|=3+4k^2 \).
若 \( k^2\ge 1 \)，则左边为 \( 4k^2-4 \)，代入上式得 \( -4=3 \)，矛盾. 所以 \( k^2<1 \). 于是 \( 4\left( 1-k^2 \right)=3+4k^2 \)，解得 \( 8k^2=1 \)，所以 \( k=\pm \frac{\sqrt{2}}{4} \).
\end{solution}

\begin{problem}
设函数 \( f(x)=x^2+mx-4\e^{nx-1} \)（\( n\in \Z \)），曲线 \( y=f(x) \) 在点 \( \left( -1,f(-1) \right) \) 处的切线方程为 \( 4x+\e^2y+\e^2+8=0 \).

\begin{enumerate}
\item 求 \( m \)，\( n \) 的值；
\item 求 \( f(x) \) 的极值点个数；
\item 求 \( y=kx-1 \)（\( k>0 \)）与 \( f(x) \) 交点个数.
\end{enumerate}

\end{problem}

\begin{solution}
（1）由切线方程 \( 4x+\e^2y+\e^2+8=0 \)，得切线斜率为 \( -\frac{4}{\e^2} \)，且切点坐标满足 \( f\left( -1 \right)=-1-\frac{4}{\e^2} \).
又 \( f\left( -1 \right)=1-m-4\e^{-n-1} \)，所以 \( 1-m-4\e^{-n-1}=-1-\frac{4}{\e^2} \)，即 \( m=2-4\e^{-n-1}+\frac{4}{\e^2} \).
再由 \( f'(x)=2x+m-4n\e^{nx-1} \)，得 \( f'\left( -1 \right)=-2+m-4n\e^{-n-1}=-\frac{4}{\e^2} \). 代入上式中的 \( m \)，得 \( \left( n+1 \right)\e^{-n-1}=2\e^{-2} \). 令 \( t=n+1 \)，则 \( t\in \Z \) 且 \( t\e^{-t}=2\e^{-2} \).
由右端大于 0，知 \( t>0 \).

当 \( t=1 \) 时，\( t\e^{-t}=\e^{-1}>2\e^{-2} \). 当 \( t=2 \) 时，\( t\e^{-t}=2\e^{-2} \).

当 \( t\ge 3 \) 时，设 \( \phi(x)=x\e^{-x} \)，则 \( \phi'(x)=\e^{-x}\left( 1-x \right)<0 \)
当 \( x\ge 3 \) 时成立.
所以 \( \phi(x) \) 在 \( \left[ 3,+\infty \right) \) 上严格递减，从而 \( t\e^{-t}\le 3\e^{-3}<2\e^{-2} \).
故唯一可能是 \( t=2 \). 所以 \( n=1 \). 代回得 \( m=2 \).

（2）由（1）知 \( f(x)=x^2+2x-4\e^{x-1} \)，于是 \( f'(x)=2x+2-4\e^{x-1} \). 设 \( g(x)=2x+2-4\e^{x-1} \)，则 \( g'(x)=2-4\e^{x-1} \). 当 \( x=1-\ln 2 \) 时，\( g'(x)=0 \). 所以 \( g(x) \) 先增后减，且 \( g\left( 1-\ln 2 \right)=2-2\ln 2>0 \).
又 \( g(-1)=-4\e^{-2}<0 \)，\( g(0)=2-4\e^{-1}>0 \)，且 \( g(1)=0 \). 结合 \( g(x) \) 的单调性可知，\( g(x)=0 \) 有 2 个不同实根，即 \( f'(x)=0 \) 有 2 个不同实根. 所以 \( f(x) \) 有 2 个极值点.

（3）令 \( F_k(x)=f(x)-\left( kx-1 \right)=x^2+\left( 2-k \right)x+1-4\e^{x-1} \)，
则直线 \( y=kx-1 \) 与曲线 \( y=f(x) \) 的交点个数，就是方程 \( F_k(x)=0 \) 的实根个数.
先看导数，\( F_k'(x)=2x+2-k-4\e^{x-1}=g(x)-k \).
由（2）知 \( g(x) \) 的最大值为 \( 2-2\ln 2 \).

若 \( k\ge 2-2\ln 2 \)，则对任意 \( x \) 都有 \( F_k'(x)\le 0 \). 当 \( k>2-2\ln 2 \) 时，\( F_k'(x)<0 \)；当 \( k=2-2\ln 2 \) 时，\( F_k'(x) \) 只在 \( x=1-\ln 2 \) 处取 0. 所以 \( F_k(x) \) 严格递减. 又 \( F_k(-2)=1+2k-4\e^{-3}>0 \)，\( F_k(0)=1-4\e^{-1}<0 \)，故此时 \( F_k(x)=0 \) 只有 1 个实根.

若 \( 0<k<2-2\ln 2 \)，则 \( F_k'(x)=0 \) 有两个实根，设它们为 \( x_1<x_2 \). 因为 \( F_k'(x)=g(x)-k \)，而 \( g(x) \) 先增后减，所以 \( x_1 \) 是极小点，\( x_2 \) 是极大点. 又 \( g(-2)=-2-4\e^{-3}<0<k \)，所以 \( x_1>-2 \).
先证 \( x_2\in \left( 1-\ln 2,1 \right) \). 由 \( x_2>1-\ln 2 \) 是显然的. 又 \( F_k'(1)=4-k-4=-k<0 \)，而在 \( x=1-\ln 2 \) 处，\( F_k'\left( 1-\ln 2 \right)=2-2\ln 2-k>0 \)，
所以由连续性知 \( x_2\in \left( 1-\ln 2,1 \right) \).
在驻点处有 \( 4\e^{x-1}=2x+2-k \). 于是对任一驻点 \( x \)，都有 \( F_k\left( x \right)=x^2+\left( 2-k \right)x+1-4\e^{x-1}=1-x^2+4\e^{x-1}\left( x-1 \right) \).
设 \( H(x)=1-x^2+4\e^{x-1}\left( x-1 \right) \)，则 \( H'(x)=-2x+4x\e^{x-1}=2x\left( 2\e^{x-1}-1 \right) \)，所以 \( H(x) \) 在 \( \left[ 1-\ln 2,1 \right] \) 上递增. 又 \( H(1)=0 \)，故 \( H\left( x_2 \right)<H(1)=0 \)，从而 \( F_k\left( x_2 \right)=H\left( x_2 \right)<0 \).
这说明 \( F_k(x) \) 的局部极大值仍小于 0. 当 \( x\ge x_1 \) 时，\( F_k(x) \le F_k\left( x_2 \right)<0 \)，没有实根. 又 \( F_k(-2)=1+2k-4\e^{-3}>0 \)，且 \( F_k(x) \) 在 \( x<x_1 \) 时严格递减，所以 \( F_k(x)=0 \) 在 \( \left( -2,x_1 \right) \) 内有且仅有 1 个实根.

综上，对任意 \( k>0 \)，直线 \( y=kx-1 \) 与 \( y=f(x) \) 都只有 1 个交点.
\end{solution}

\begin{problem}
设 \( A=\left( a_{ij} \right)_{m\times n} \) 是一个 \( m \) 行 \( n \) 列的数阵，且数阵中的每一项 \( a_{ij}\in \left\{ -1,1 \right\} \)，都等于 1 或 \(-1\). 若对任意 \(1\le i\)，\(\ p\le m\)，\(1\le j\)，\(\ q\le n\)，其中 \( i\neq p \)，\( j\neq q \)，且 \( \left| \left| i-p \right|-\left| j-q \right| \right|=2 \)，都有 \( a_{ij}+a_{iq}+a_{pj}+a_{pq}=0 \)，则称数阵 \(A\) 具有性质 \(P\).

\begin{enumerate}
\item 判断下列两个数表是否具有性质 \(P\)：

\begin{center}
\renewcommand{\arraystretch}{1.2}
\begin{tabular}{cc}
\(A_1=\)
\begin{tabular}{|c|c|c|c|}
\hline
1 & 1 & \(-1\) & 1 \\
\hline
\(-1\) & 1 & \(-1\) & 1 \\
\hline
\end{tabular}
&
\(A_2=\)
\begin{tabular}{|c|c|c|c|}
\hline
1 & 1 & \(-1\) & \(-1\) \\
\hline
1 & \(-1\) & 1 & \(-1\) \\
\hline
\(-1\) & 1 & \(-1\) & 1 \\
\hline
\end{tabular}
\end{tabular}
\end{center}
\item 在所有具有性质 \(P\) 的 \( 5\times 4 \) 数阵中，1 的个数最多是多少？
\item 若 \( m=n=6 \)，\( a_{11}=1 \)，且数阵 \(A\) 具有性质 \(P\)，证明：对任意 \( i \)，\( j\in \left\{ 1,2,3,4,5,6 \right\} \)，都有 \( a_{ij}=a_{i1}a_{1j} \).
\end{enumerate}

\end{problem}

\begin{solution}
先证一个后面反复要用的结论：若 \(u\)，\(v\)，\(s\)，\(t\in \left\{ -1,1 \right\}\)，且 \( u+v+s+t=0 \)，则 \( uv=st \). 理由是 \( u+v=-(s+t) \)，两边平方得 \( u^2+2uv+v^2=s^2+2st+t^2 \)，即 \( uv=st \). 因而对任意满足性质 \(P\) 的四个角，都有 \( a_{ij}a_{iq}=a_{pj}a_{pq} \).

（1）对 \(A_1\)，只有两行，所以 \( \left| i-p \right|=1 \). 又要满足 \( \left| \left| i-p \right|-\left| j-q \right| \right|=2 \)，只能取 \( \left| j-q \right|=3 \)，即只需检查第 1 列与第 4 列. 于是 \( a_{11}+a_{14}+a_{21}+a_{24}=1+1-1+1=2\ne 0 \)，
所以 \(A_1\) 不具有性质 \(P\).

对 \(A_2\)，当 \( \left| i-p \right|=2 \) 时，需要 \( \left| j-q \right|=0 \) 或 4，但题设要求 \( j\ne q \)，且 4 列数表中 \( \left| j-q \right| \le 3 \)，因此这种情形不存在. 所以只需检查相邻两行与第 1，4 列：\( a_{11}+a_{14}+a_{21}+a_{24}=1-1+1-1=0 \)，且 \( a_{21}+a_{24}+a_{31}+a_{34}=1-1-1+1=0 \).
故 \(A_2\) 具有性质 \(P\).

（2）设 \( x_i=a_{i1}+a_{i4} \)，\( y_i=a_{i2}+a_{i3} \)（\(i=1\)，\(2\)，\(3\)，\(4\)，\(5\)）.
则每个 \(x_i\)，\(y_i\in \left\{ -2,0,2 \right\}\).
当两行相邻时，必须取第 1，4 列，所以 \( x_i+x_{i+1}=0 \)（\(i=1\)，\(2\)，\(3\)，\(4\)）. 于是 \( x_1=-x_2=x_3=-x_4=x_5 \)，故第 1，4 列中 1 的个数 \( N_{14}=5+\frac{x_1+x_2+x_3+x_4+x_5}{2}=5+\frac{x_1}{2}\le 6 \).

当两行相距 3 时，必须取相邻两列. 由第 2，3 列可得 \( y_1+y_4=0 \) 且 \( y_2+y_5=0 \)，所以第 2，3 列中 1 的个数 \( N_{23}=5+\frac{y_1+y_2+y_3+y_4+y_5}{2}=5+\frac{y_3}{2}\le 6 \).
因此整个 \( 5\times 4 \) 数阵中，1 的个数至多为 \( N_{14}+N_{23}\le 6+6=12 \).

下面给出一个能取到 12 的例子：
\[
\begin{array}{|c|c|c|c|}
\hline
1 & -1 & -1 & 1 \\
\hline
-1 & 1 & 1 & -1 \\
\hline
1 & 1 & 1 & 1 \\
\hline
-1 & 1 & 1 & -1 \\
\hline
1 & -1 & -1 & 1 \\
\hline
\end{array}
\]
它其中共有 12 个 1. 又在 \( 5\times 4 \) 数阵中，满足题设的情形只可能是“行距 1，列距 3”或“行距 3，列距 1”两类. 对这个例子，\( \left( 1 \right)+\left( 1 \right)+\left( -1 \right)+\left( -1 \right)=0 \)，且 \( \left( 1 \right)+\left( -1 \right)+\left( -1 \right)+\left( 1 \right)=0 \)，
对应的其余同类四角和也都相同为 0，所以它满足性质 \(P\). 因此 1 的个数最多是 12.

（3）先证第 4 行与第 1 行成比例. 因为 \( \left| 4-1 \right|=3 \)，对 \(j=1\)，\(2\)，\(3\)，\(4\)，\(5\) 都有 \( a_{1j}+a_{1,j+1}+a_{4j}+a_{4,j+1}=0 \). 由前面的乘积结论，\( a_{1j}a_{1,j+1}=a_{4j}a_{4,j+1} \)，从而 \( a_{4j}a_{1j}=a_{4,j+1}a_{1,j+1} \). 所以 \( a_{4j}a_{1j} \) 与 \( j \) 无关. 又 \( a_{11}=1 \)，故 \( a_{4j}=a_{41}a_{1j} \)（\(j=1\)，\(2\)，\(\cdots\)，\(6\)）.

同理，\( a_{5j}a_{2j}=a_{51}a_{21} \)（\(j=1\)，\(2\)，\(\cdots\)，\(6\)），并且 \( a_{6j}a_{3j}=a_{61}a_{31} \)（\(j=1\)，\(2\)，\(\cdots\)，\(6\)）.

下面证明第 2 行也与第 1 行成比例.

由 \( \left| 2-1 \right|=1 \)，取第 1，4 列，得 \( a_{21}a_{24}=a_{11}a_{14}=a_{14} \)，故 \( a_{24}=a_{21}a_{14} \).
由 \( \left| 4-2 \right|=2 \)，取第 1，5 列，得 \( a_{21}a_{25}=a_{41}a_{45} \). 又 \( a_{45}=a_{41}a_{15} \)，且 \( a_{41}^2=1 \)，所以 \( a_{21}a_{25}=a_{15} \)，故 \( a_{25}=a_{21}a_{15} \).
再由 \( \left| 2-1 \right|=1 \)，取第 2，5 列，得 \( a_{22}a_{25}=a_{12}a_{15} \). 代入 \( a_{25}=a_{21}a_{15} \)，得 \( a_{22}=a_{21}a_{12} \).
由 \( \left| 4-2 \right|=2 \)，取第 2，6 列，得 \( a_{22}a_{26}=a_{42}a_{46} \). 而 \( a_{42}=a_{41}a_{12} \)，\( a_{46}=a_{41}a_{16} \)，所以 \( a_{22}a_{26}=a_{12}a_{16} \). 代入 \( a_{22}=a_{21}a_{12} \)，得 \( a_{26}=a_{21}a_{16} \).
再由 \( \left| 2-1 \right|=1 \)，取第 3，6 列，得 \( a_{23}a_{26}=a_{13}a_{16} \). 代入 \( a_{26}=a_{21}a_{16} \)，得 \( a_{23}=a_{21}a_{13} \).
综上，\( a_{2j}=a_{21}a_{1j} \)（\(j=1\)，\(2\)，\(\cdots\)，\(6\)）.
下面证明第 3 行也与第 1 行成比例.
由 \( \left| 5-3 \right|=2 \)，取第 1，5 列，得 \( a_{31}a_{35}=a_{51}a_{55} \). 由上面的关系 \( a_{5j}a_{2j}=a_{51}a_{21} \) 和第 2 行结论，可得 \( a_{5j}=a_{51}a_{21}a_{2j}=a_{51}a_{1j} \)（\(j=1\)，\(2\)，\(\cdots\)，\(6\)）. 于是 \( a_{55}=a_{51}a_{15} \)，从而 \( a_{31}a_{35}=a_{15} \)，故 \( a_{35}=a_{31}a_{15} \).
由 \( \left| 3-2 \right|=1 \)，取第 2，5 列，得 \( a_{22}a_{25}=a_{32}a_{35} \). 而 \( a_{22}=a_{21}a_{12} \)，\( a_{25}=a_{21}a_{15} \)，故 \( a_{32}a_{35}=a_{12}a_{15} \). 代入 \( a_{35}=a_{31}a_{15} \)，得 \( a_{32}=a_{31}a_{12} \).
由 \( \left| 3-1 \right|=2 \)，取第 2，6 列，得 \( a_{32}a_{36}=a_{12}a_{16} \). 代入 \( a_{32}=a_{31}a_{12} \)，得 \( a_{36}=a_{31}a_{16} \).
由 \( \left| 3-2 \right|=1 \)，取第 3，6 列，得 \( a_{23}a_{26}=a_{33}a_{36} \). 而 \( a_{23}=a_{21}a_{13} \)，\( a_{26}=a_{21}a_{16} \)，故 \( a_{33}a_{36}=a_{13}a_{16} \). 代入 \( a_{36}=a_{31}a_{16} \)，得 \( a_{33}=a_{31}a_{13} \).
由 \( \left| 3-4 \right|=1 \)，取第 1，4 列，得 \( a_{31}a_{34}=a_{41}a_{44} \). 又 \( a_{44}=a_{41}a_{14} \)，所以 \( a_{31}a_{34}=a_{14} \)，故 \( a_{34}=a_{31}a_{14} \).
再由 \( \left| 3-4 \right|=1 \)，取第 2，5 列，得 \( a_{32}a_{35}=a_{42}a_{45} \).
因为 \( a_{42}=a_{41}a_{12} \)，\( a_{45}=a_{41}a_{15} \)，所以上式右边等于 \( a_{12}a_{15} \)，与前面一致.
由 \( \left| 3-4 \right|=1 \)，取第 3，6 列，得 \( a_{33}a_{36}=a_{43}a_{46} \).
因为 \( a_{43}=a_{41}a_{13} \)，\( a_{46}=a_{41}a_{16} \)，所以上式右边等于 \( a_{13}a_{16} \)，与前面一致.

因此 \( a_{3j}=a_{31}a_{1j} \) 对 \(j=1\)，\(2\)，\(\cdots\)，\(6\) 都成立.

于是第 5 行满足 \( a_{5j}=a_{51}a_{21}a_{2j}=a_{51}a_{1j} \)，
这里用到了 \( a_{5j}a_{2j}=a_{51}a_{21} \) 以及 \( a_{2j}^2=1 \).
第 6 行满足 \( a_{6j}=a_{61}a_{31}a_{3j}=a_{61}a_{1j} \)，
这里用到了 \( a_{6j}a_{3j}=a_{61}a_{31} \) 以及 \( a_{3j}^2=1 \).
再结合第 1，2，3，4 行的结论，可知对任意 \( i\in \left\{ 1,2,3,4,5,6 \right\} \)，都有 \( a_{ij}=a_{i1}a_{1j} \)（\(j=1\)，\(2\)，\(\cdots\)，\(6\)）.
证毕.
\end{solution}
